\chapter{Introduction}
\label{c:introduction}

\noindent
Emotion recognition in text tries to identify the affective state expressed or implied in written
language. Sentiment analysis typically reduces to positive and negative polarity; emotion
classification instead assigns finer-grained categories such as joy, anger, fear, shame, pride or
sadness. The task is hard because emotional meaning is rarely stated outright. It arises
from the interaction of linguistic content, situational context, and the assumptions a reader
makes about the person speaking.

The major challenge is ambiguity. Some texts allow a variety of reasonable emotional readings,
due to underspecified context, an emotionally complicated situation, or the relative position of
writer and reader. Even short, seemingly simple statements can be read in competing ways.
Consider:

\begin{quote}
I can't believe I actually did that.
\end{quote}

\noindent
Without further context, this could be a statement of pride, embarrassment, regret, excitement,
or disbelief. Background assumptions tend to resolve such ambiguity for the reader, and one of
the most accessible is an assumption about \emph{who} is speaking. Attributed to someone running
their first marathon, the sentence invites one emotional reading; attributed to someone
recounting a lapse of judgment, it invites a different one, yet the words are the same.

\section{Motivation}
\label{intro:motivation}

\noindent
The thesis has two motivations.

The first is computational. Large language models are now widely used for text classification,
emotion recognition included, and are known to be sensitive to prompt framing and contextual
cues. If minimal author information systematically changes a model's emotional interpretation,
then an emotion classifier's output is shaped by the social context supplied in the prompt, not
by text content alone. That bears directly on how such systems are evaluated and deployed.

The second is cognitive and social: do models reflect human interpretive behaviour when dealing
with ambiguous emotional content? Humans demonstrably use assumptions about a speaker to
disambiguate. Whether models do the same, and to what extent and in what direction, is an
empirical question. It has not been directly investigated by manipulating perceived authorship
while holding the text constant.

\section{Problem Statement}
\label{intro:problem}

\noindent
Current research on emotion recognition largely treats disagreement among annotators as
noise to be removed. Datasets often include a single aggregated label for each text, most
frequently by majority vote, and models are trained and evaluated against that label. Built into
this practice is an assumption that fails for a large portion of ambiguous texts: that there is
one correct emotional reading and that annotators who do not stay with that reading are simply
in error.

The problem addressed here has two parts. The first is whether perceived author identity has a
systematic effect on emotion interpretation, in both humans and models, or only affects it
incidentally. The second is whether the normal practice of collapsing judgments to a single label
can detect such an influence at all. Suppose author framing causes a shift in the
\emph{distribution} of plausible readings without necessarily changing which reading is most
popular. An analysis based on a majority label will then record no change where an actual change
occurred.

\section{Research Approach}
\label{intro:approach}

\noindent
The thesis manipulates perceived authorship experimentally and keeps the textual content fixed.
Three hundred texts from the crowd-enVent corpus \cite{Troiano2023CrowdEnVent} were stratified
into three groups: unambiguous control texts, author-independent ambiguous texts, and
author-relevant ambiguous texts. Each text was framed neutrally and under two contrasting author
personas. Human annotations were collected for all framings, and seven large language models
labelled every text across four framings and six prompt settings.

What distinguishes this work is its choice of representation. All judgments, human and model, are
modelled as full probability distributions over the label set of eleven emotions rather than as a
single label. Framing effects are then shifts in that distribution: their \emph{magnitude} is the total variation
distance, their \emph{direction} the signed change in probability mass across emotions. Using this,
one can distinguish a text that was only slightly reinterpreted from one that was substantially
reinterpreted but happened to retain the same most-popular answer.

\section{Contributions}
\label{intro:contributions}

\noindent
The contributions of this thesis are as follows.

\begin{itemize}
    \item An experimental design that dissociates the influence of perceived authorship on
    emotion interpretation by changing author framing while holding text fixed, implemented in
    parallel on human annotators and seven large language models.

    \item A distributional analysis framework in which each judgment is an eleven-dimensional
    probability vector, with the magnitude and direction of framing-induced shift measured
    separately and never mixed together.

    \item Empirical evidence that author framing measurably shifts both human and model emotion
    judgments, together with a control condition that dissociates the effect of the persona's
    content from the mere presence of an author attribution.

    \item A direct measure of what majority-label analysis gives up: the fraction of judgments
    that shifted but would have been recorded as unchanged.

    \item An explicit account of what these data can and cannot establish, including a
    sensitivity analysis showing that, at this sample size, the predicted ordering across
    ambiguity categories cannot be separated from baseline disagreement.
\end{itemize}

\section{Thesis Structure}
\label{intro:structure}

\noindent
Chapter~\ref{c:literature} examines research on emotion recognition and annotator disagreement,
and on framing effects arising from perceived authorship, and identifies the gap this thesis
addresses. Chapter~\ref{c:methodology} describes the corpus, how texts were categorised by
ambiguity type, and the author-framing manipulation, together with the research questions and
hypotheses. The entire pipeline is recorded in chapter~\ref{c:implementation}: prompt
construction, model evaluation, the human annotation study, the generation of emotion probability
vectors, the metrics derived from them, and the quality-control and statistical procedures.
Chapter~\ref{c:result} is a hypothesis-by-hypothesis report of the results. These results are
explained in relation to the research questions in chapter~\ref{c:discussion}, and the
limitations of the design are set out, along with recommendations for future research.
Chapter~\ref{c:conclusion} concludes.
